\section*{\texorpdfstring{\S\CurrentExerciseGroup}{Section \CurrentExerciseGroup}}
\phantomsection
\addcontentsline{toc}{section}{\texorpdfstring{Exercises for \S\CurrentExerciseGroup}{Exercises for Section \CurrentExerciseGroup}}

\begin{enumerate}
\def\labelenumi{\arabic{enumi})}
\tightlist
\item
  Let \(\mathbf{T}\) and \(\mathbf{X}\) be two locally compact spaces, \(\mu\) a positive measure on \(\mathbf{T}\) , \(\pi\) a continuous and \(\mu\) -proper mapping of \(\mathbf{T}\) into \(\mathbf{X}\) , and let \(\nu = \pi(\mu)\) .
\end{enumerate}

\begin{enumerate}
\def\labelenumi{\alph{enumi})}
\item
  Show that for every numerical function \(g \geqslant 0\) , defined and lower semi-continuous on \(X\) , \(g \circ \pi\) is lower semi-continuous on \(T\) and \(\nu^{*}(g) = \mu^{*}(g \circ \pi)\) .
\item
  Show that if \(\mathbf{f}\) is a \(\nu\) -integrable function with values in a Banach space, then \(\mathbf{f} \circ \pi\) is \(\mu\) -integrable and \(\int \mathbf{f} d\nu = \int (\mathbf{f} \circ \pi) d\mu\) . Is the converse valid without a supplementary condition (cf.~§4, Exer. 2b)?
\end{enumerate}

¶ 2) Let T and X be two locally compact spaces, P the set of proper continuous mappings of T into X, equipped with the topology of compact convergence. Let H be a subset of P such that, for every compact subset K of X, the union of the sets \(\overline{\pi}^{-1}(K)\) , where \(\pi\) runs over H, is relatively compact in T.

\begin{enumerate}
\def\labelenumi{\alph{enumi})}
\item
  Show that the mapping \((\pi, \lambda) \mapsto \pi(\lambda)\) of \(\mathrm{H} \times \mathcal{M}(\mathrm{T})\) into \(\mathcal{M}(\mathrm{X})\) is not necessarily continuous when each of the spaces \(\mathcal{M}(\mathrm{T})\) and \(\mathcal{M}(\mathrm{X})\) is equipped with the quasi-strong topology (Ch. III, §1, Exer. 8).
\item
  If B is a bounded subset of \(\mathcal{M}(\mathrm{T})\) , show that the restriction to \(\mathrm{H} \times \mathrm{B}\) of the mapping \((\pi, \lambda) \mapsto \pi(\lambda)\) is continuous when each of the spaces \(\mathcal{M}(\mathrm{T}), \mathcal{M}(\mathrm{X})\) is equipped with the vague topology (make use of Prop. 15 of Ch. III, §1).
\item
  Show that the mapping \((\pi, \lambda) \mapsto \pi(\lambda)\) of \(\mathrm{H} \times \mathcal{M}_{+}(\mathrm{T})\) into \(\mathcal{M}_{+}(\mathrm{X})\) is continuous when \(\mathcal{M}(\mathrm{T})\) and \(\mathcal{M}(\mathrm{X})\) are equipped with the vague topology (cf.~Ch. III, §1, Exer. 10). d) Give an example where \(\mathrm{T} = \mathrm{X}\) is compact and the mapping \((\pi, \lambda) \mapsto \pi(\lambda)\) of \(\mathcal{P} \times \mathcal{M}(\mathrm{T})\) into \(\mathcal{M}(\mathrm{X}) = \mathcal{M}(\mathrm{T})\) is not continuous, when \(\mathcal{M}(\mathrm{T})\) is equipped with the vague topology (take for \(\mathrm{T}\) the torus \(\mathbf{T}\) and use Exer. 2 b) of Ch. III, §4).
\end{enumerate}

¶ 3) Let T be a noncompact locally compact space, \(\mu\) a positive measure on T, and \(\mathcal{G}\) the group of homeomorphisms of T. Give an example showing that the mapping \(\pi \mapsto \pi(\mu)\) of \(\mathcal{G}\) into \(\mathcal{M}(T)\) is not necessarily continuous when \(\mathcal{G}\) is equipped with the topology of compact convergence and \(\mathcal{M}(T)\) with the vague topology (take for T the subspace of R formed by 0 and the points n and 1/n (n an integer \(\geqslant 1\) ) and for \(\mu\) the measure defined by \(\mu(\{0\}) = 0\) , \(\mu(\{1/n\}) = 1/n^2\) , \(\mu(\{n\}) = 1\) ).

\begin{enumerate}
\def\labelenumi{\arabic{enumi})}
\setcounter{enumi}{3}
\item
  For every locally compact space T, let \(\mathcal{M}^{c}(T)\) be the subspace of \(\mathcal{M}(T)\) formed by the real measures with compact support. If T and X are two locally compact spaces, and \(\pi\) is a continuous mapping of T into X, then \(\pi\) is \(\lambda\) -proper and \(\pi(\lambda) \in \mathcal{M}^{c}(X)\) for every measure \(\lambda \in \mathcal{M}^{c}(T)\) . Give an example showing that the mapping \(\lambda \mapsto \pi(\lambda)\) of \(\mathcal{M}^{c}(T)\) into \(\mathcal{M}^{c}(X)\) is not necessarily continuous when \(\mathcal{M}^{c}(T)\) and \(\mathcal{M}^{c}(X)\) are equipped with the vague topology (or with the quasi-strong topology (Ch. III, §1, Exer. 8)) and that \(\pi\) is not a proper mapping (cf.~Ch. III, §2, Exer. 1).
\item
  \begin{enumerate}
  \def\labelenumii{\alph{enumii})}
  \tightlist
  \item
    Let \(\rho\) and \(\sigma\) be two positive measures on \(\mathbf{R}\) . Show that if \(\rho(]a,b]) = \sigma(]a,b])\) for every semi-open interval \(]a,b]\) , then \(\rho = \sigma\) .
  \end{enumerate}
\end{enumerate}

\begin{enumerate}
\def\labelenumi{\alph{enumi})}
\setcounter{enumi}{1}
\tightlist
\item
  Let I be an interval of \(\mathbf{R}\) , with left end-point \(\alpha\) and right end-point \(\beta\) , and let \(\psi\) be an increasing finite numerical function defined on I; extend \(\psi\) to \(\mathbf{R}\) in any way whatsoever. In order that \(\psi\) be proper for the measure \(\varphi_{\mathrm{I}} \cdot \mu\) ( \(\mu\) Lebesgue measure on \(\mathbf{R}\) ), it is necessary and sufficient that the following two conditions hold: \(1^{\circ}\) either \(\alpha\) is finite, or \(\alpha = -\infty\) and \(\lim_{x \to -\infty} \psi(x) = -\infty\) ; \(2^{\circ}\) either \(\beta\) is finite, or \(\beta = +\infty\) and
\end{enumerate}

\(\lim_{x\to +\infty}\psi (x) = +\infty .\)

Suppose these conditions verified. For every \(y \in \mathbf{R}\) , let \(\theta(y)\) be the supremum of the set of \(x \in I\) such that \(\psi(x) \leqslant y\) (the supremum being equal to \(\alpha\) if this set is empty). Show that \(\theta\) is a finite numerical function, increasing and right-continuous in \(\mathbf{R}\) . If \(\nu\) is the image of \(\varphi_{\mathrm{I}} \cdot \mu\) under \(\psi\) , then \(\nu(]a,b]) = \theta(b) - \theta(a)\) for every semi-open interval \(]a,b]\) in \(\mathbf{R}\) .

\begin{enumerate}
\def\labelenumi{\alph{enumi})}
\setcounter{enumi}{2}
\tightlist
\item
  Conversely, show that for every finite numerical function \(\theta\) , increasing and right-continuous in \(\mathbf{R}\) , there exists one and only one positive measure \(\nu\) on \(\mathbf{R}\) such that \(\nu(]a,b]) = \theta(b) - \theta(a)\) for every semi-open interval \(]a,b]\) of \(\mathbf{R}\) (the 'Stieltjes measure on \(\mathbf{R}\) defined by \(\theta'\) ; one writes \(\int f d\theta\) instead of \(\int f d\nu\) ); moreover, every positive measure on \(\mathbf{R}\) may be obtained in this way. Under what condition is the measure \(\nu\) diffuse? What is then the image of \(\nu\) under \(\theta\) ?
\end{enumerate}

\begin{enumerate}
\def\labelenumi{\arabic{enumi})}
\setcounter{enumi}{5}
\tightlist
\item
  \begin{enumerate}
  \def\labelenumii{\alph{enumii})}
  \tightlist
  \item
    Let K be the Cantor triadic set in the interval {[}0,1{]} (GT, IV, §2, No.~5). Show that there exist a diffuse positive measure \(\nu\) on \(\mathbf{R}\) , with support K and total mass 1, and a continuous increasing function \(\theta\) , defined on \(\mathbf{R}\) , such that \(\theta(\mathbf{R}) = \mathbf{I}\) and \(\theta(\nu) = \varphi_{\mathrm{I}} \cdot \mu\) ( \(\mu\) the Lebesgue measure; cf.~GT, IV, §8, Exer. 16).
  \end{enumerate}
\end{enumerate}

\begin{enumerate}
\def\labelenumi{\alph{enumi})}
\setcounter{enumi}{1}
\tightlist
\item
  Deduce from a) that there exists a diffuse positive measure on R, alien to Lebesgue measure, whose support is all of I (in each interval J contiguous to K and contained in I, take a measure proportional to the image of \(\nu\) under an affine mapping of I onto J; then proceed by recursion).
\end{enumerate}

\begin{enumerate}
\def\labelenumi{\arabic{enumi})}
\setcounter{enumi}{6}
\tightlist
\item
  Let \(\mu\) be Lebesgue measure on \(\mathbf{R}\) , I the interval \([0,1]\) of \(\mathbf{R}\) . If one sets \(\theta(x) = |x|\) , \(\theta\) is proper for the measure \(\varphi_{\mathrm{I}} \cdot \mu\) and one has \(\theta(\varphi_{\mathrm{I}} \cdot \mu) = \varphi_{\mathrm{I}} \cdot \mu\) . Give an example of a set negligible for \(\varphi_{\mathrm{I}} \cdot \mu\) , whose image under \(\theta\) is not measurable for \(\varphi_{\mathrm{I}} \cdot \mu\) (cf.~Ch. IV, §4, Exer. 8).
\end{enumerate}

¶ 8) Let T be a compact space, μ a diffuse positive measure on T of total mass 1.

\begin{enumerate}
\def\labelenumi{\alph{enumi})}
\tightlist
\item
  Show that there exists a continuous mapping \(\pi\) of T onto E = {[}0,1{]} such that the image of \(\mu\) under \(\pi\) is the Lebesgue measure \(\lambda\) on E. (Proceed as in the proof of Urysohn's theorem (GT, IX, §4, Th. 1) by defining, for every \(t\) such that \(0 \leqslant t \leqslant 1\) , an open set \(\mathrm{U}(t) \subset \mathrm{T}\) , quadrable for \(\mu\) (Ch. IV, §5, Exer. 17), such that \(\mathrm{U}(0) = \varnothing\) , \(\mathrm{U}(1) = \mathrm{T}\) , \(\overline{\mathrm{U}(t)} \subset \mathrm{U}(t')\) for \(t < t'\) , and finally \(\mu(\mathrm{U}(t)) = t\) ; one uses Exer. 7 of §5 to prove that if V, W are two quadrable open sets in T such that \(\overline{\mathrm{V}} \subset \mathrm{W}\) and \(\mu(\mathrm{V}) < \mu(\mathrm{W})\) , then there exists a quadrable open set U such that \(\overline{\mathrm{V}} \subset \mathrm{U} \subset \overline{\mathrm{U}} \subset \mathrm{W}\) and such that
\end{enumerate}

\[
\frac {1}{3} \mu (\mathrm{W-V}) \leqslant \mu (\mathrm{U-V}) \leqslant \frac {2}{3} \mu (\mathrm{W-V}).
\]

Finally, one applies Exer. 5 a).)

\begin{enumerate}
\def\labelenumi{\alph{enumi})}
\setcounter{enumi}{1}
\item
  Deduce from a) that there exist subsets of T that are not \(\mu\) -measurable (Ch. IV, §4, Exer. 8), that there exist sequences in \(\mathcal{L}^{p}(\mathrm{T},\mu)\) that convergence in mean of order p to 0 but do not converge to 0 at any point of T, for \(1 \leqslant p < +\infty\) (Ch. IV, §3, Exer. 1), and sequences \((f_{n})\) such that \((\widetilde{f}_{n})\) converges to 0 for the weakened topology of \(\mathrm{L}^{p}(\mathrm{T},\mu)\) but does not converge in measure to 0 (§5, Exer. 19).
\item
  Suppose, moreover, that T is metrizable. Show that there exist a \(\mu\) -negligible subset N of T, a \(\lambda\) -negligible subset M of E, and a homeomorphism \(\pi\) of E-M onto T-N such that, if one extends (arbitrarily) \(\pi\) to a mapping \(\varphi\) of E into T, and \(\pi^{-1}\) to a mapping \(\psi\) of T into E, then \(\varphi(\lambda)=\mu\) and \(\psi(\mu)=\lambda\) . (Using Exer. 17 of Ch. IV, §5, show that for every integer n\textgreater0, there exists a finite partition of T formed by a \(\mu\) -negligible set and by quadrable open sets, of diameter \(\leqslant1/n\) (for a metric compatible with the topology of T) and of measure \(\leqslant1/n\) . Proceeding recursively, deduce therefrom, by passage to the limit, the existence of a continuous mapping f of E-D into T, where D is a countable subset of E, such that \(f(\lambda)=\mu\) ; show that one can arrange that f be a homeomorphism of E-D onto a subset of T of measure 1 (cf.~§8, Exer. 14).)
\end{enumerate}

¶ 9) Let \(\mu\) and \(\nu\) be two diffuse positive measures on a locally compact space T. Show that, for \(\nu\) to be a measure with base \(\mu\) , it suffices that every \(\mu\) -measurable subset of T be also \(\nu\) -measurable (argue by contradiction using Th. 3 and Prop. 13 of §5, and Exer. 8 b) of §6).

\begin{enumerate}
\def\labelenumi{\arabic{enumi})}
\setcounter{enumi}{9}
\tightlist
\item
  In the interval \(\mathbf{I} = ]0,1[\) of \(\mathbf{R}\) , define the function \(g\) by \(g(t) = -\sqrt{n}\) for \(\frac{1}{2}\left(\frac{1}{n} + \frac{1}{n+1}\right) < t \leqslant \frac{1}{n}\) and \(g(t) = \sqrt{n}\) for \(\frac{1}{n+1} < t \leqslant \frac{1}{2}\left(\frac{1}{n} + \frac{1}{n+1}\right) (n = 1,2,\ldots)\) , and extend \(g\) by 0 in \(\mathbf{R} - \mathbf{I}\) ; \(g\) is integrable for Lebesgue measure. Set
\end{enumerate}

\[
\mathrm{G} (x) = \int_ {0} ^ {x} g (t) d t;
\]

show that the function \(f\) , equal to \(1 / \sqrt{x}\) on \([-1, +1]\) and to 0 elsewhere, which is integrable for Lebesgue measure, is such that \(t \mapsto f(\mathrm{G}(t))g(t)\) is not integrable on I.

¶ 11) The notations are those of No.~5; g is a locally μ-integrable numerical function on I.

\begin{enumerate}
\def\labelenumi{\alph{enumi})}
\item
  For G to be a \(\lambda\) -proper mapping of I into G(I), it is necessary and sufficient that the following conditions be fulfilled: \(1^{\circ}\) the limits G(a+) and G(b−) exist in \(\overline{\mathbf{R}}\) ; \(2^{\circ}\) if G(a+) ∈ G(I) then g is \(\mu\) -integrable on the interval {]} \(a, x_0[\) ; \(3^{\circ}\) if G(b−) ∈ G(I) then g is \(\mu\) -integrable on the interval {]} \(x_0, b[\) .
\item
  Suppose that the limits \(\mathrm{G}(a+)\) and \(\mathrm{G}(b-)\) exist in \(\overline{R}\) . Show that if f is such that \(t \mapsto \mathbf{f}(\mathrm{G}(t))g(t)\) is integrable for Lebesgue measure on I, then f is integrable on the interval with end-points \(\mathrm{G}(a+)\) and \(\mathrm{G}(b-)\) for Lebesgue measure, and formula (13) holds. (Reduce to the case that \(x_{0}\) is one of the numbers a, b; if, for example, \(x_{0} = a\) , observe that there exists in I a strictly increasing sequence \((b_{n})\) tending to b, such that the sequence \((\mathrm{G}(b_{n}))\) is either increasing or decreasing; apply Prop. 4 of Ch. IV, §4, No.~3 and Lebesgue's theorem).
\end{enumerate}

¶ 12) a) Let g be a finite numerical function defined on a compact interval I = {[}a, b{]} of R. One calls right-expansion set (resp. left-expansion set) of g in I the set of \(x \in I\) such that there exists \(y \in I\) for which x < y (resp. x \textgreater{} y) and \(g(x) < g(y)\) . Show that if g is continuous, then the right- (resp. left-) expansion set of g in I is open in I and that, if {]} \(\alpha, \beta[\) is a connected component \(^{1}\) of this set, then \(g(\alpha) = g(\beta)\) and \(g(x) \leqslant g(\alpha)\) for \(\alpha < x < \beta\) .

\begin{enumerate}
\def\labelenumi{\alph{enumi})}
\setcounter{enumi}{1}
\item
  Let \(r_1, r_2\) be two real numbers such that \(0 \leqslant r_1 < r_2\) . Let \(g\) be an increasing continuous function on I; let \(E'\) be the left-expansion set of \(g(x) - r_1x\) in I, \(E''\) the union of the right-expansion sets of \(g(x) - r_2x\) in each of the intervals that are the closures of the connected components of \(E'\) . If \(\mu\) is Lebesgue measure on I, show, using a), that \(\mu(E'') \leqslant \frac{r_1}{r_2} \mu(E')\) .
\item
  One calls the upper and lower right derivates at a point \(x \in I\) , of a finite numerical function \(f\) defined on \(I\) , the respective numbers
\end{enumerate}

\[
\begin{array}{l} \mathrm{D} _ {r} ^ {+} f (x) = \operatorname * {l i m s u p} _ {h \to 0,   h > 0} \left(f (x + h) - f (x)\right) / h \\ \mathrm{D} _ {r} ^ {-} f (x) = \operatorname * {l i m i n f} _ {h \to 0,   h > 0} \left(f (x + h) - f (x)\right) / h. \end{array}
\]

Similarly, one calls the upper and lower left derivates of \(f\) at the point \(x\) , the respective numbers

\[
\begin{array}{l}\mathrm{D} _ {l} ^ {+} f (x) = \underset {h \rightarrow 0, h <   0} {\lim \sup} \left(f (x + h) - f (x)\right) / h\\\mathrm{D} _ {l} ^ {-} f (x) = \underset {h \rightarrow 0, h <   0} {\lim \inf} \left(f (x + h) - f (x)\right) / h.\end{array}
\]

Show that if \(f\) is continuous and increasing, then the set of \(x \in I\) where \(\mathrm{D}_r^+ f(x) = +\infty\) and the set of \(x \in I\) where \(\mathrm{D}_l^- f(x) < \mathrm{D}_r^+ f(x)\) have measure zero for \(\mu\) (for every rational number \(r > 0\) (resp. for every pair \((r_1, r_2)\) of rational numbers such that \(0 \leqslant r_1 < r_2\) ), consider the set of \(x \in I\) such that \(\mathrm{D}_r^+ f(x) > r\) (resp. \(\mathrm{D}_l^- f(x) < r_1\) and \(\mathrm{D}_r^+ f(x) > r_2\) simultaneously), and apply \(b\) )). From this, deduce that for almost every \(x \in I\) , \(f\) admits a finite derivative (`Lebesgue's theorem') (apply the preceding result also to \(-f(-x)\) ).

¶ 13) In a compact interval I = {[}a, b{]} of R, let \((f_n)\) be a sequence of continuous, increasing finite functions, such that \(f_n(a) = 0\) for all n and such that the sum \(s(x) = \sum_{n=1}^{\infty} f_n(x)\) is finite on I. Show that there exists a set N ⊂ I, negligible for Lebesgue measure, such that, for every \(x \in I - N\) , the derivatives \(f_n'(x)\) and \(s'(x)\) exist and \(s'(x) = \sum_{n=1}^{\infty} f_n'(x)\) (`Fubini's theorem'). (Use Exer. 12. Observe that the series with general term \(f_n'(x)\) is convergent almost everywhere; setting

\[
s _ {n} (x) = \sum_ {k = 1} ^ {n} f _ {k} (x),
\]

consider next a subsequence \((s_{n_k})\) such that the series with general term \(s(x) - s_{n_k}(x)\) is convergent on I, and apply the preceding remark to this series.)

\begin{enumerate}
\def\labelenumi{\arabic{enumi})}
\setcounter{enumi}{13}
\item
  Let \(f\) be a numerical function defined on a compact interval \(I = [a, b]\) of \(\mathbf{R}\) , integrable on \(I\) for Lebesgue measure. If one sets \(F(x) = \int_{a}^{x} f(t) dt\) , show that \(F\) admits almost everywhere in \(I\) a derivative equal to \(f\) (reduce to the case that \(f \geqslant 0\) ; then, using Exer. 13, consider successively the case that \(f\) is lower semi-continuous then the general case (cf.~Ch. IV, §4, No.~4, Cor. of Th. 3)).
\item
  Denoting by \(\mu\) the Lebesgue measure on \(\mathbf{R}\) , a point \(x \in \mathbf{R}\) is said to be a density point of a subset A of \(\mathbf{R}\) if, when \(h\) and \(k\) tend to 0 through values \(>0\) , the quotient \(\mu^{*}(\mathrm{A} \cap [x - h, x + k]) / (h + k)\) tends to 1. Show that the set of points of an arbitrary subset A of \(\mathbf{R}\) that are not density points of A is negligible for \(\mu\) . (Reduce to the case that A is contained in a compact interval \([a, b]\) , and consider the function \(s_{\mathrm{A}}(x) = \mu^{*}(\mathrm{A} \cap [a, x])\) . Take a decreasing sequence \((\mathrm{A}_n)\) of open sets containing A, such that the series with general term \(s_{\mathrm{A}_n}(x) - s_{\mathrm{A}}(x)\) is convergent on \([a, b]\) , and use Exer. 13.)
\item
  \begin{enumerate}
  \def\labelenumii{\alph{enumii})}
  \tightlist
  \item
    Let \(g\) be a function integrable for Lebesgue measure on \(\mathbf{E} = [0,1]\) . Set \(\mathrm{G}(x) = \int_0^x g(t)dt\) ; show that at every point \(x \in \mathbf{E}\) where \(g\) is continuous, \(\mathbf{G}\) admits a derivative equal to \(g\) .
  \end{enumerate}
\end{enumerate}

\begin{enumerate}
\def\labelenumi{\alph{enumi})}
\setcounter{enumi}{1}
\item
  Give an example of a function \(g\) that is bounded and is continuous almost everywhere in \(\mathbf{E}\) , such that \(\mathbf{G}\) has no right derivative at the points of an uncountable subset of \(\mathbf{E}\) (cf.~FRV, I, §2, Exer. 8).
\item
  Show that the function equal to \(x^2 \sin \frac{1}{x^2}\) for \(x \neq 0\) , and to 0 for \(x = 0\) , admits a derivative at every point of \(E\) , but this derivative is not integrable for Lebesgue measure on \(E\) .
\end{enumerate}

\begin{enumerate}
\def\labelenumi{\arabic{enumi})}
\setcounter{enumi}{16}
\item
  Let T and X be two locally compact spaces, \(\pi\) a continuous mapping of T into X. Show that if \(\pi\) is \(\mu\) -proper for every positive measure \(\mu\) on T, then \(\pi\) is proper. (Argue by contradiction, using Prop. 7 of GT, I, §10.)
\item
  In Prop. 4 b), the conclusion may fail if one omits the hypothesis that \(\pi'\) is continuous. (Take \(T = R\) , \(T' = \overline{R}\) , \(T'' = R\) ; take for \(\mu\) Lebesgue measure, for \(\pi\) the canonical injection; set \(\pi'(x) = x\) for \(x \in R\) and \(\pi'(\pm \infty) = 0\) .)
\item
  Let \(\mathrm{T}, \mathrm{X}, \mathrm{Y}\) be locally compact spaces, \(\pi : \mathrm{T} \to \mathrm{X}\) , \(\pi' : \mathrm{X} \to \mathrm{Y}\) mappings, \(\pi'' = \pi' \circ \pi\) , \(\mu\) a real measure on \(\mathrm{T}\) . It can happen that \(\pi\) is \(\mu\) -proper and \(\pi'\) is \(\pi(\mu)\) -proper, without \(\pi''\) being \(\mu\) -proper. (Take \(\mathrm{T} = \mathbf{R} \times \{0, 1\}\) , \(\mathrm{X} = \mathbf{R}\) , \(\mathrm{Y} = \{0\}\) , \(\pi = \mathrm{pr}_1\) ; take for \(\mu\) the measure that induces Lebesgue measure (resp. the opposite of Lebesgue measure) on \(\mathbf{R} \times \{0\}\) (resp. \(\mathbf{R} \times \{1\}\) ) canonically identified with \(\mathbf{R}\) .)
\item
  Let \(\mathrm{T},\mathrm{X},\mathrm{Y}\) be three locally compact spaces, \(\mu\) a positive measure on \(\mathrm{T}\) , \(t\mapsto \lambda_t\) a \(\mu\) -adequate mapping of \(\mathrm{T}\) into \(\mathcal{M}_{+}(\mathrm{X})\) , \(\nu = \int \lambda_t d\mu (t)\) , and \(\pi\) a \(\nu\) -proper mapping of \(\mathrm{X}\) into \(\mathrm{Y}\) . Assume one of the following hypotheses: a) \(\mathrm{Y}\) is countable at infinity and \(\nu\) is moderated; b) \(\mathrm{Y}\) is countable at infinity and \(\lambda_t\) is bounded locally almost everywhere in \(\mathrm{T}\) . Then \(\pi\) is \(\lambda_t\) -proper locally almost everywhere in \(\mathrm{T}\) , the mapping \(t\mapsto \pi (\lambda_t)\) is \(\mu\) -adequate, and its integral is equal to \(\pi (\mu)\) .
\item
  Let \(\mathrm{T},\mathrm{X},\mathrm{Y}\) be three locally compact spaces, \(\pi : \mathrm{T} \to \mathrm{X}\) , \(\pi': \mathrm{X} \to \mathrm{Y}\) universally measurable mappings. Then \(\pi' \circ \pi\) is universally measurable. (Let \(\mu\) be a positive measure on \(\mathrm{T}\) . To prove that \(\pi' \circ \pi\) is \(\mu\) -measurable, reduce to the case that \(\mathrm{T}\) is compact and \(\pi\) is continuous, and use the fact that \(\pi'\) is measurable for \(\pi(\mu)\) .)
\item
  Let \(X, Y\) be two compact spaces, \(U\) a continuous linear mapping of the Banach space \(\mathcal{C}(X; \mathbf{R})\) into the Banach space \(\mathcal{C}(Y; \mathbf{R})\) , \(\pi : X \to Y\) a continuous mapping. Assume that for every \(y \in Y\) and every function \(f \in \mathcal{C}(X; \mathbf{R})\) ,
\end{enumerate}

\[
\inf _ {\pi (x) = y} f (x) \leqslant (U \cdot f) (y) \leqslant \sup _ {\pi (x) = y} f (x).
\]

Show that under these conditions there exists a vaguely continuous mapping \(\sigma: y \mapsto \sigma_y\) of Y into \(\mathcal{M}_+(X)\) such that, for every \(y \in Y\) , \(\sigma_y\) is carried by \(\overline{\pi}^{-1}(y)\) and one has \((U \cdot f)(\pi(x)) = \langle f, \sigma_{\pi(x)} \rangle\) for all \(x \in X\) ; converse. (Consider the transpose \(^t U\) .)

\setcounter{exercisegroup}{7}
\edef\CurrentExerciseGroup{\arabic{exercisegroup}}
